\section{Introduction}
\label{sec:introduction}

Advanced Persistent Threats (APTs) remain among the most damaging classes of cyber attacks, often persisting undetected in enterprise networks for months~\cite{provdetector,hrat}.
To counter these threats, provenance-based Host Intrusion Detection Systems (HIDS) reconstruct system-level provenance graphs from audit logs and apply machine learning to identify malicious behavior~\cite{unicorn,shadewatcher,threatrace}.
These systems have demonstrated strong detection capabilities: by modeling the causal relationships between processes, files, and network connections, they can surface attack campaigns that evade traditional signature-based defenses.

However, the effectiveness of provenance-based HIDS has created an adversarial arms race.
Recent work has explored \emph{mimicry attacks}---techniques that inject benign-looking operations into provenance graphs to disguise malicious behavior~\cite{provninja,malguise}.
ProvNinja~\cite{provninja}, the state of the art, selects camouflage operations based on a \emph{regularity score} that measures how frequently an operation appears in normal system behavior.
The underlying assumption is intuitive: operations that are common in benign environments should be effective at masking attacks.

\paragraph{The functionality problem.}
This assumption is flawed.
Consider a Linux host environment where \texttt{kill -9 \textlangle pid\textrangle} is a highly regular operation.
Injecting this command during a firefox backdoor attack destroys the attacker's foothold, terminating the compromised process and all child processes within it.
The operation is common, but it is catastrophic to attack functionality.
More broadly, any camouflage operation that modifies, deletes, or disrupts resources the attack depends on can silently break the attack chain---regardless of how ``normal'' it appears.
We find that ProvNinja's regularity-based selection suffers from exactly this problem: its camouflage not only fails to improve evasion but actually \emph{backfires}, increasing ML detection rates from 94\% to 100\% (\S\ref{sec:evaluation}).

\paragraph{The stealthiness problem.}
Even when camouflage operations do not break the attack, na\"ive insertion introduces detectable artifacts.
Modern detection systems employ not only ML-based classifiers but also rule-based statistical checks: edge-count thresholds ($>\mu + 3\sigma$), KL divergence bounds on operation distributions, and time-window burst detectors~\cite{unicorn,shadewatcher}.
Blindly adding edges to the provenance graph---even benign ones---inflates edge counts, skews the distribution of operation types, and creates temporal bursts that trip these secondary defenses.
Prior mimicry attacks target only the ML detection layer and ignore these rule-based checks entirely.

\paragraph{Comparison with other domains.}
Adversarial evasion has been studied for Android malware, where BagAmmo~\cite{bagammo} wraps injected API calls in try-catch blocks to neutralize side effects and EvadeDroid~\cite{evadedroid} isolates added components from the original app logic.
These isolation mechanisms do not transfer to provenance graphs: every inserted operation \emph{actually executes} on the host, producing real system calls, modifying real files, and creating real network connections.
There is no try-catch wrapper for \texttt{rm -rf} at the syscall level.
The provenance graph setting therefore demands a fundamentally different approach to ensuring that camouflage does not compromise attack functionality.

\paragraph{Our approach.}
We present \textbf{SafeMimic}, a dependency-aware adversarial mimicry attack that formally separates the \emph{safety} of camouflage operations from their \emph{stealthiness}.
SafeMimic first extracts the attack's dependency set $\Sdep(A)$---the processes, files, and network connections that the attack chain requires---and computes the effect set $\Seffect(b)$ of each candidate camouflage operation $b$.
It then constructs a \emph{safe pool} $\Bsafe = \{b \mid \Seffect(b) \cap \Sdep(A) = \emptyset\}$, guaranteeing that no selected operation can interfere with attack execution.
Only after safety is established does SafeMimic optimize for stealthiness, selecting operations from $\Bsafe$ via $\DKL$ minimization to match the target environment's normal distribution across both operation types and temporal patterns.

Figure~\ref{fig:hero} illustrates the contrast between regularity-based selection (ProvNinja) and SafeMimic's dependency-aware approach.
While ProvNinja may select high-regularity operations that destroy attack state, SafeMimic provably avoids such operations and simultaneously produces distribution-matched camouflage that evades both ML classifiers and rule-based checks.

\begin{figure*}[t]
    \centering
    \fbox{\parbox{0.95\textwidth}{\centering\vspace{2em}\textbf{[Hero Figure Placeholder]}\\[0.5em]
    (a) ProvNinja: regularity-based selection $\to$ unsafe ops break attack + excess edges trigger rules\\
    (b) SafeMimic: dependency-aware filtering $\to$ safe pool + KL-matched insertion\\
    (c) Summary: attack success rate + ML detection rate comparison\vspace{2em}}}
    \caption{Overview of SafeMimic vs.\ ProvNinja. (a)~ProvNinja selects camouflage by regularity without safety checks---high-regularity operations like \texttt{kill -9 \textlangle pid\textrangle} can destroy attack state, and excess edges create detectable frequency anomalies. (b)~SafeMimic first filters via dependency analysis ($\Bsafe$), then selects via distribution matching to evade both ML and rule-based detection. (c)~Only SafeMimic achieves simultaneous attack preservation (100\%) and ML evasion (0\% detection).}
    \label{fig:hero}
\end{figure*}

\paragraph{Contributions.}
We make the following contributions:

\begin{enumerate}
    \item \textbf{Formal dependency framework.} We define the dependency set $\Sdep(A)$, effect set $\Seffect(b)$, and safe operation pool $\Bsafe$, and prove a Safety Theorem guaranteeing that operations drawn from $\Bsafe$ cannot compromise attack functionality (\S\ref{sec:design}).

    \item \textbf{Regularity $\neq$ safety.} We prove (Proposition~\ref{prop:regularity}) that an operation's regularity score is neither necessary nor sufficient for safety, providing a direct theoretical critique of ProvNinja's design (\S\ref{sec:formulation}).

    \item \textbf{Dynamic safe space.} We introduce $\Bsafe(t)$, a time-varying safe pool that evolves as the attack progresses through its stages, accommodating the changing dependency structure of multi-step attack campaigns (\S\ref{sec:design}).

    \item \textbf{Distribution-matched insertion.} We design a dual-layer evasion strategy that selects operations from $\Bsafe$ to minimize $\DKL$ divergence from the normal distribution while respecting edge-count, distribution, and temporal constraints imposed by rule-based detectors (\S\ref{sec:design}).

    \item \textbf{Comprehensive evaluation.} We evaluate SafeMimic against three state-of-the-art HIDS (Unicorn~\cite{unicorn}, ShadeWatcher~\cite{shadewatcher}, ThreaTrace~\cite{threatrace}) using four DARPA TC attack scenarios and three DARPA Transparent Computing datasets. SafeMimic achieves 100\% attack success with 0\% ML detection, reducing anomaly scores by 95\% on average (\S\ref{sec:evaluation}).
\end{enumerate}

\noindent
Our evaluation reveals a striking finding: ProvNinja's camouflage \emph{backfires}.
On attacks where the original (uncamouflaged) detection rate is 94\%, ProvNinja's inserted operations push it to 100\%.
Meanwhile, random camouflage insertion preserves attack success only 24--67\% of the time.
SafeMimic is the first mimicry attack that provably guarantees both attack preservation and evasion effectiveness against multi-layer detection.
